
% 来源：EPUB part169.xhtml
\section{腾讯控股}

\subsection{通过社交媒体将流量货币化了}

\BookQuestion{网友：每次你分析公司总能直达本质，一针见血。估摸着大道要出手腾讯了。}

长长的坡，厚厚的雪！（2018-10-15）

最近确实已经投了一些了，但比例还很小。用苹果换似乎有点不舍得，我似乎更懂苹果些。要再仔细想想。（2018-10-16）

我觉得10年后回过头来看，腾讯赚的钱（我买的份额占的比例）应该比我存银行的利息要多。（2019-03-16）

\BookQuestion{网友：您去年建仓腾讯肯定是基于多年的观察，您投资前有做过什么功课吗？}

我投资的标准很简单：商业模式，企业文化，合理价钱。所以当商业模式和企业文化不错的公司价钱也不错时，我就会有兴趣。（2019-03-17）

\BookQuestion{网友：未来十年可能会创造2000-3000亿元的年利润，请问这个数字是怎么来的？看了您以前的文章您几乎都能提前五六年预测苹果和茅台的收入和利润，而且它们俩还真做到了！}

其实都是瞎猜的。生意模式好，企业文化好，基本上你就可以简单地希望他们会正常增长，所以给个五年十年的，有增长不奇怪啊。关键是要看懂生意和商业模式。（2019-03-17）

\BookQuestion{网友：你对苹果说得多，对腾讯很少提，你如何看腾讯的企业文化？}

我对苹果了解比较多，对腾讯了解相对少。直觉上觉得腾讯商业模式非常不错，企业文化也还可以哈。（2019-04-11）

\BookQuestion{网友：您最认可腾讯哪些方面？}

马化腾人不错，而且年轻，微信及微信支付的影响非常大且深远。我们确实买了一点腾讯，但比例很小。（2019-10-15）

\BookQuestion{网友：您觉得腾讯10年后会怎样？能看到50年吗？}

简单讲，腾讯就是通过建立社交媒体，将流量货币化了。感觉10年后腾讯应该依然强大，但50年就不好说了。我对腾讯的理解度远不如苹果，所以投的比例也小很多。（2020-10-09）

\BookQuestion{网友：大道评价一下阿里跟腾讯吧？}

我觉得阿里的护城河没有腾讯的宽，但今日的出现让人觉得腾讯的护城河也没那么宽了。投资确实很难哈。（2022-03-09）

\BookQuestion{网友：感觉大道在腾讯的买入上还是要保守些，相对苹果来说。}

是的，目前还不到1\%，还不算实质性投入。还在想……（2022-08-02）

\BookQuestion{网友：确定的好是第一位的，从长期来看只有确定的好才可能是便宜的，或许苹果便宜些。}

是的，腾讯对我来说确定性确实比苹果小不少，这也是现阶段一直下不了大决心多买的原因。苹果下来的时候我完全不用下决心就能下买单哈。不过，腾讯对我而言依然是非卖品，找到一个商业模式好的公司不容易。我猜未来几年内我还会有机会进一步加仓腾讯的，目前确实还是有些事情没想清楚。（2022-08-07）

我对苹果确实完全不操心，掉多少都不往心里去。腾讯总是觉得懂得不透，下不了重手。

\BookQuestion{网友：最疑虑的点是哪些？}

未来现金流看不透。（2022-10-05）

\BookQuestion{网友：大道唯独对茅台和苹果买入的价格没那么敏感，基本是有钱就买。}

是的，苹果跌得多的时候我一定会买的，茅台经常会，腾讯有时候也会买一点。（2022-10-30）

% 来源：EPUB part170.xhtml
\subsection{用户为本，科技向善}

腾讯目前的股数比2016年还多，所以这些年的回购实际上全部都变成了员工股？苹果过去10年的股数已经减少了40\%左右了。

苹果在过去10年里每一年的股票都是减少的。腾讯除了2023年比2022年股数略有减少外，其他年都是增加的，而且2023年的股数居然比2016年还多。这种回购实际上就是在给员工发奖金哈。

\BookQuestion{网友：这或许就是大道一直不重仓腾讯的缘由吧？}

原因之一吧。我们买的到底是什么很重要。（2023-12-26）

我一直在投资腾讯啊，只是投得比较少而已。有些顾虑挥之不去，只好慢慢来了。（2023-12-28）

\BookQuestion{网友：请教下你如何看腾讯的企业愿景？}

以前是“做一家受人尊敬的互联网公司”，现在是“用户为本，科技向善”。我非常喜欢新的愿景！太牛了！八个字，该有的都有了！老的那个我很久以前就说过，不是很喜欢。

\BookQuestion{网友：“应该有的都有了”，指的应该是消费者导向和利润之上的追求这两条。同样很喜欢，现在的愿景，英文是Value for Users, Tech for}

Good。

我喜欢你的解释！（2023-12-12）

% 来源：EPUB part171.xhtml
\subsection{根据自己的机会成本做决定}

\BookQuestion{网友：未来10到15年，您认为腾讯能到的市值是多少？}

市值是别人愿意出多少钱买这家公司的意思。我不知道别人愿意出多少钱买这家公司，所以我不知道腾讯能到多少市值。不过，我认为腾讯这家公司大概率10到15年内一年可以赚2000-3000亿人民币，愿意出多少钱买这家公司完全取决于你自己的机会成本。（2019-03-14）

如果考虑腾讯的投资部分，P/E实际上已经低过20了。（2024-08-15）

\BookQuestion{网友：老段这个算法怎么回事？腾讯投资部分公允5000多亿，账面价值3000多亿，看利润用的是哪个，扣非净利润，自由现金流？}

我没算过，但你可以假设腾讯把投资都卖掉，把钱还完债务后全部用来回购，完了你再算P/E就比较真实了。这是假设，没必要但可以这么做。（2024-08-17）

拼多多还是一家初创公司，未来变数还很多。腾讯相比之下要更成熟。该说的都说过了哈。（2024-10-28）

\BookQuestion{网友：我觉得腾讯的模式比苹果还牛。}

让腾讯上6000亿美元试试？腾讯确实是不错的公司，但要成为苹果还需要很多东西，包括大环境的变化。顺便说一下什么叫生意模式：生意模式就是赚钱的模式，最后赚得越多的模式就应该越好。苹果现在一年的利润比中国所有互联网公司（包括腾讯）从开办以来赚的利润的总和都要多，你说腾讯的模式更好的逻辑是啥？（2013-03-04）

腾讯早晚会再上6000亿的。（2024-10-19）
